\subsection{Fixed affine rank at every temperature}
\label{sec:fixedrankalltemp}

Polynomial prefill can remove the numerical temperature dependence
for every fixed affine rank.  The price is a much larger index than
the near-linear index for one affine line.  The construction first
indexes the order of all projected keys, then uses the decreasing-weight
envelope of Lemma~\ref{lem:newmonotoneenvelope}.

Point location in arrangements of hyperplanes is classical.  For
example, Meiser~\cite{meiser1993pointlocation} gives polynomial space
at each fixed dimension and logarithmic query dependence on the
number of hyperplanes.  We give a slower elementary construction
here so that its preprocessing, rational arithmetic, and treatment
of tied scores are explicit.  No geometric oracle is assumed.

\begin{lemma}[An elementary index for hyperplane arrangements]
\label{lem:elementaryarrangement}
Let $r\ge1$, and let $\mathcal H$ be a family of at most $H$
proper affine hyperplanes in $\mathbb R^r$, given by rational
coefficients of encoding length at most $b$.  There is a finite
family of nonempty, open, full-dimensional convex cells such that:
\begin{enumerate}
 \item every cell avoids all hyperplanes in $\mathcal H$, and every
 input affine form has a constant strict sign on each cell;
 \item the cells have disjoint interiors, and their closures cover
 $\mathbb R^r$;
 \item their number is at most $(H+1)^{2^r-1}$;
 \item each cell has a stored rational representative; a rational
 query point can be located in the closure of one stored cell using
 at most $(2^r-1)\lceil\log_2(H+1)\rceil+O(r)$ comparisons of
 rational affine forms.
\end{enumerate}
For each fixed $r$, all coefficients and representatives have
encoding length $O_r(b+1)$.  The preprocessing work and stored bits
are at most
\[
 (H+1)^{2^r-1}\operatorname{poly}_r(b+\log(H+2)),
\]
and query bit work is
$\operatorname{poly}_r(b+b_q+\log(H+2))$, where $b_q$ bounds the
query-coordinate encoding lengths.  Hyperplanes with identical zero
sets, empty zero sets, and identically zero forms can be removed by
exact rational tests.
\end{lemma}

\paragraph{Proof idea.}
Project the hyperplanes one coordinate at a time.  Their crossing sets partition the lower-dimensional space into regions with a fixed vertical order.  Lifting those regions gives cells and a recursive binary-search index that also handles boundary ties.

\begin{proof}
In one dimension, sort the distinct rational zeros and store the
open intervals between consecutive zeros, including the two
unbounded intervals.  Their representatives may be chosen as
midpoints, or one unit beyond an outer endpoint.  With no zeros
use the representative zero.  Binary search locates a query in
the closure of one interval.  At a zero, either adjacent interval
is valid.  There are at most $H+1$ intervals, establishing the base
case.

Suppose $r\ge2$, write a point as $(x,y)$ with
$x\in\mathbb R^{r-1}$, and separate the affine forms according to
whether their coefficient of $y$ is zero.  A form with nonzero
coefficient has zero set
\[
 y=f_i(x),
\]
where $f_i$ is rational affine.  Remove identical functions among
the $f_i$, and let their number be $K\le H$.  In the projection
space, collect the nontrivial zero sets of all differences
$f_i-f_j$ and of all input forms independent of $y$.  This gives
at most
\[
 H'\le \binom K2+H-K\le (H+1)^2-1
\]
proper affine hyperplanes.  Constant nonzero forms have empty zero
sets and are discarded; zero forms add no constraint.

Construct the lower-dimensional index recursively.  On each of
its cells $C$, the distinct functions $f_i$ have a strict order
that is independent of $x\in C$.  Obtain this order by evaluating
them at the stored representative $x_C$ and sorting.  Write the
order as
$f_{i_1}<\cdots<f_{i_K}$ on $C$.  Above $C$, create the $K+1$
open cells
\[
 \{(x,y):x\in C,\ f_{i_j}(x)<y<f_{i_{j+1}}(x)\},
 \qquad 0\le j\le K,
\]
with the usual endpoints $f_{i_0}=-\infty$ and
$f_{i_{K+1}}=+\infty$.  When $K=0$, use $C\times\mathbb R$.
Every such cell is nonempty, open, and convex.  A representative
uses $x_C$ and the midpoint of the two finite threshold values;
at an unbounded end take the finite threshold minus or plus one.
All original affine forms have a fixed strict sign in a lifted
cell, including those independent of $y$.

For a query $(x,y)$, first recurse to obtain a cell $C$ whose
closure contains $x$.  Its stored threshold order remains a weak
order at $x$, by continuity.  A binary search through this order
finds an index $j$ for which
\[
 f_{i_j}(x)\le y\le f_{i_{j+1}}(x).
\]
Equal threshold values cause no difficulty.  We now verify that
the returned lifted cell really has $(x,y)$ in its closure.
Choose $x_m\in C$ tending to $x$.  Its corresponding open vertical
interval is nonempty.  The distance of $y$ from this interval
tends to zero because each finite endpoint is a continuous affine
function and $y$ lies between the limiting endpoints.  Choose a
point $y_m$ in the interval whose distance from $y$ is at most
that interval distance plus $1/m$.  Then
$(x_m,y_m)\to(x,y)$.  This proves the closure claim even when the
two limiting endpoints coincide.  Disjointness of cell interiors
follows from the recursive projection and the disjoint vertical
intervals.

Let $N_r(H)$ bound the number of cells.  The construction gives
\[
 N_r(H)\le(H+1)N_{r-1}(H')
 \le(H+1)\bigl((H+1)^2\bigr)^{2^{r-1}-1}
 =(H+1)^{2^r-1}.
\]
At successive projection levels, the logarithm of the number of
hyperplanes plus one grows by at most a factor of two.  Summing
the binary-search costs gives the stated comparison count.

The projected coefficients are obtained from two previous affine
forms by rational division and subtraction.  The number of such
projection stages is at most $r-1$.  A representative at each
stage is obtained by substituting the preceding representative
into at most two affine forms and taking a midpoint or adding one.
Each operation combines only $O(r)$ rational numbers.  Induction
therefore bounds every coefficient and representative by
$C_r(b+1)$ bits for a constant $C_r$ independent of $H$.

At each projection cell, sorting its $K$ threshold values costs
$O(K\log(K+2))$ rational comparisons.  The stored threshold order
has $K$ indices, while that projection cell gives $K+1$ lifted
cells.  Thus this storage and sorting cost are bounded by a
logarithmic factor times the number of cells created, with rational
arithmetic charged.  Generating all pairs of forms costs $O(H^2)$
operations at a projection step, which is absorbed by the displayed
bound whenever its dimension is at least two.  The one-dimensional
base step uses direct sorting and creates no pairwise projection.
This proves the preprocessing and storage bounds.  Query tests
use only the precomputed rational coefficients and supplied query
coordinates, establishing the bit bound.
\end{proof}

\begin{theorem}[An exact index for any fixed affine rank]
\label{thm:fixedrankalltemperature}
Fix $r\ge0$.  Let $T\ge2$ cached dyadic keys in $\mathbb Q^n$
have affine dimension at most $r$, with dyadic scalar values in
$[-1,1]$.  Suppose cache coordinates, values, and a fresh explicit dyadic
query have at most $b$ bits each.  Let the nonnegative temperature
$\beta$ be dyadic of at most $b$ bits, or let $\beta=\sqrt n$.
For
\[
 B=16+b+\lceil\log_2(T+n+2)\rceil,
\]
there is an exact attention-index sampler with the following bounds:
\[
 \begin{split}
 \text{prefill work and stored bits}
 &\le T^{\max\{1,2^r-1\}}\operatorname{poly}_r(n,B),\\
 \text{expected online bit work}
 &=O_r\bigl((n+\log T)B^4\bigr).
 \end{split}
\]
Its expected number of proposed keys is at most $17/8$.
The bounds have no dependence on the numerical size of $\beta$
or of the scores beyond their input encoding lengths.  A fresh
value sign gives the exact attention-coordinate mean.  Composing
with the unit tanh row and the bounded two-logit head changes the
online bound by only an absolute factor.
\end{theorem}

\paragraph{Proof idea.}
The query matters only through its projections onto the affine key span.  Index every possible projected order in fixed dimension, then apply the decreasing-weight envelope to the order selected by the query.

\begin{proof}
Find the actual affine rank $r'\le r$ by exact rational Gaussian
elimination.  If $r'=0$, all scores agree and a uniform index
suffices.  Otherwise choose one key $c$ and $r'$ independent key
differences $u_1,\ldots,u_{r'}$.  Express each key uniquely as
\[
 k_j=c+\sum_{a=1}^{r'}\xi_{j,a}u_a.
\]
One nonsingular $r'\times r'$ coordinate minor determines each
$\xi_j$.  Cramer's rule, or rational elimination, shows that its
numerators and denominators have $O_r(b+1)$ bits.  The basis
selection, coefficient construction, and rank-promise verification
cost $T\operatorname{poly}_r(n,B)$ bit operations.

For a query $q$, define its $r'$ projected coordinates
\[
 v_a=\langle q,u_a\rangle.
\]
The score of key $j$ is a common additive term plus
$(\beta/n)\langle v,\xi_j\rangle$.  Thus its order is determined
by the signs of the at most $\binom{T}{2}$ forms
\[
 \langle v,\xi_j-\xi_k\rangle.
\]
Their homogeneity saves one dimension in the arrangement index.
For each coordinate $a\le r'$ and sign $\sigma\in\{-1,1\}$,
construct the chart $v_a=\sigma$.  It has $r'-1$ free coordinates
$y_b$, $b\ne a$, and its pairwise forms are
\[
 \sigma(\xi_{j,a}-\xi_{k,a})
 +\sum_{b\ne a}(\xi_{j,b}-\xi_{k,b})y_b.
\]
Discard constant forms and apply
Lemma~\ref{lem:elementaryarrangement} to the proper hyperplanes
in each chart.  When $r'=1$, each chart is a single point and no
point-location structure is needed.  At the representative of
each chart cell, sort the $T$ projected keys and store their
permutation.  This costs $T\log(T+2)$ comparisons and $T$ indices
per cell.  For $r'\ge2$, the total size of the $2r'$ chart indices
is at most
\[
 2r'T(H+1)^{2^{r'-1}-1}\operatorname{poly}_r(n,B)
 \le T^{2^{r'}-1}\operatorname{poly}_r(n,B),
\]
where $H\le T^2$ and factors depending only on $r$ are absorbed.
For $r'=1$ this is $T\operatorname{poly}(n,B)$.  Together with
the rank-zero case these are bounded by the stated exponent.

At query time compute $v$.  If $v=0$, all scores agree and use a
uniform index.  Otherwise select an index $a$ with $v_a\ne0$, put
$\sigma=\operatorname{sign}(v_a)$, and query its chart at
$y_b=v_b/|v_a|$.  Multiplication by the positive number $|v_a|$
recovers the original pairwise score differences.  Locate a chart
cell whose closure contains $y$.  Its stored descending score order
is weakly valid at $y$: each pairwise difference has a fixed sign
on the open cell and has the corresponding weak sign on its closure.
The discarded constant forms have the same sign everywhere, and
identically zero forms represent ties everywhere.  Thus the stored
order is weakly valid for the original query, including every tie.
The order is unchanged by any positive temperature; at temperature
zero use a uniform index.

Apply the dyadic decreasing-weight envelope from the proof of
Theorem~\ref{thm:newrankone}.  It needs only the sorted permutation,
$O(\log T)$ score evaluations, and a constant expected number of
proposed-member score evaluations.  There is no need to read the
entire permutation online.  Its absolute exponential enclosures,
dyadic proposal floors, and final acceptance comparisons give
exactness and the bound $17/8$ on the number of proposals.

The projection uses $nr'$ dyadic products and additions.  Its
coordinates have $O_r(B)$ bits.  Chart selection and rational
division preserve that encoding bound.  Point location uses
$O_r(\log T)$ rational affine comparisons with
$O_r(B)$-bit coefficients.  These costs fit in the displayed
$O_r((n+\log T)B^4)$ bound.  The static index has polynomial size
at fixed $r$, so its addresses require only $O_r(B)$ bits.
After the query projection is available, evaluate each shifted score
using the stored affine coordinates:
\[
 a_j-a_{\max}=\frac\beta n
          \langle v,\xi_j-\xi_{\max}\rangle.
\]
This takes $O_r(1)$ rational operations on $O_r(B)$-bit numbers.
The result is rational, or a rational divided by $\sqrt n$ under
standard scaling, exactly the two cases of the finite-exponential
routine.  Thus the $O(\log T)$ block scores and the constant expected
number of member scores do not repeat the length-$n$ projection.
Known value signs and
the two scalar factories complete the claim about the final token.
\end{proof}

The construction is an existence result with an explicit polynomial
prefill budget.  The exponent grows rapidly with the affine rank.
For rank one it is near-linear in $T$, with the sharper arithmetic
bound given in Theorem~\ref{thm:newrankone}.  For rank two the
$T^3\operatorname{poly}(n,B)$ construction also has a direct angular
description: sort the $O(T^2)$ pairwise tie directions in the query
plane and store one key order per angular sector.  An exact orientation
comparison and binary search locate the query direction, and a
direction on a boundary may use either adjacent weak order.
Here each homogeneous tie line contributes its two opposite
rational rays.  Sort rays by their half-plane and the sign of their
two-by-two determinant.  After duplicate rays are removed, each
sector has angle less than $\pi$, except when there is only one
tie line.  In the former case the sum of its two boundary vectors
is a rational interior representative; in the latter case a
quarter-turn of a boundary vector represents each of the two open
half-planes.  These representatives give the stored orders.
The same half-plane and determinant comparisons locate a nonzero
query ray by binary search; the zero query has all scores tied.
With no tie lines every key projection is identical.  This gives
an explicit finite-arithmetic construction of the stated rank-two
bound.
The general fixed-rank theorem is compatible with
Theorem~\ref{thm:indexedlowerlogtwo}, whose hard affine ranks grow
with the context length.
